\section{System Integration and Command-Line Interface}

\subsection{Pipeline Runner and Orchestration}

The \texttt{run\_correlator} function orchestrates the pipeline: initialises components from configuration, iterates over data chunks processing through F-Engine and Delay Engine, correlates spectra via X-Engine, accumulates until integration completes, then writes output.
                
                if config.max_integrations and integration_count >= config.max_integrations:
                    break
        
        if config.max_integrations and integration_count >= config.max_integrations:
            break
    
    # Save configuration for reproducibility
    config.to_yaml(output_dir / 'config.yaml')
    
    return 0  # Success
\end{verbatim}

This runner encapsulates the entire processing logic in approximately 50 lines of code. Component initialisation proceeds first, validating configuration and allocating resources. The main processing loop iterates over data chunks, applies transformations sequentially, and writes outputs. Breaking from nested loops when reaching maximum integrations requires checking the condition at multiple exit points—an inelegant but straightforward approach that avoids complex control flow.

Progress logging informs users of processing status:

\begin{verbatim}
print(f"Processing chunk {chunk_idx}...")
print(f"  Saved integration {integration_count} -> visibility_{integration_count:04d}.{ext}")
print(f"Complete: {integration_count} integrations written to {output_dir}/")
\end{verbatim}

These print statements provide feedback during long-running processes. For production deployments, replacing print with proper logging (Python's \texttt{logging} module) would enable configurable verbosity and log file output, but print suffices for the research prototype.

\subsection{Command-Line Argument Parsing}

The command-line interface translates user arguments into configuration objects. Python's \texttt{argparse} module structures the interface:

\begin{verbatim}
parser = argparse.ArgumentParser(
    description="FX Correlator for Radio Telescope Arrays",
    formatter_class=argparse.ArgumentDefaultsHelpFormatter
)

parser.add_argument('--config', type=str, help='Configuration file path')
parser.add_argument('--n-ants', type=int, help='Number of antennas')
parser.add_argument('--n-channels', type=int, help='FFT size')
parser.add_argument('--sample-rate', type=float, help='Sample rate (Hz)')
parser.add_argument('--integration-time', type=float, help='Integration time (s)')
parser.add_argument('--window-type', choices=['rectangular', 'hanning', 'hamming', 'blackman'])
parser.add_argument('--mode', choices=['simulated', 'file'])
parser.add_argument('--input-file', type=str, help='Input data file')
parser.add_argument('--output-dir', type=str, help='Output directory')
parser.add_argument('--max-integrations', type=int, help='Maximum integrations')
# ... additional arguments

args = parser.parse_args()
\end{verbatim}

The parser defines expected arguments with types, choices, and help text. The \texttt{ArgumentDefaultsHelpFormatter} includes default values in help messages, documenting the baseline configuration. After parsing, the code constructs a config object by layering values:

\begin{verbatim}
# Start with defaults or load from file
if args.config:
    config = CorrelatorConfig.from_yaml(args.config)
else:
    config = CorrelatorConfig()

# Override with command-line arguments
if args.n_ants is not None:
    config.n_ants = args.n_ants
if args.n_channels is not None:
    config.n_channels = args.n_channels
# ... repeat for all parameters

# Invoke correlator
return run_correlator(config)
\end{verbatim}

This pattern implements the configuration hierarchy: defaults, then file, then arguments. The conditional assignment (\texttt{if args.param is not None}) distinguishes between "argument not provided" (None) and "argument explicitly set to some value" (non-None), preventing None from overwriting valid file-loaded values.

\subsection{Dual-Mode Interface: Development and Production}

The system entry point (\texttt{\_\_main\_\_.py}) dispatches to mode-specific handlers:

\begin{verbatim}
def main():
    parser = argparse.ArgumentParser(...)
    parser.add_argument('mode', choices=['dev', 'prod'], 
                        help='Operation mode')
    args, remaining = parser.parse_known_args()
    
    if args.mode == 'dev':
        from correlator.cli.dev import start_development_cli
        start_development_cli()
    elif args.mode == 'prod':
        from correlator.cli.prod import start_production_cli
        start_production_cli()
\end{verbatim}

Mode selection occurs before parsing the full argument set, allowing mode-specific CLIs to define different parameters. For instance, production mode might add \texttt{--stream-address} for network sources while development mode adds \texttt{--sim-duration} for simulation control. This separation keeps interface complexity manageable: users see only relevant options for their chosen mode.

The development mode CLI prioritises ease of experimentation:

\begin{verbatim}
def start_development_cli():
    parser = argparse.ArgumentParser(description="Development Mode")
    parser.add_argument('command', choices=['run', 'generate'], 
                        help='run: correlate | generate: make test data')
    
    if command == 'run':
        # Add simulation parameters
        parser.add_argument('--sim-duration', type=float, default=10.0)
        parser.add_argument('--sim-snr', type=float, default=20.0)
        # Parse and execute
        args = parser.parse_args(remaining)
        config = build_config_from_args(args)
        config.data_source = 'simulated'
        run_correlator(config)
\end{verbatim}

Production mode emphasises reliability and integration:

\begin{verbatim}
def start_production_cli():
    parser = argparse.ArgumentParser(description="Production Mode")
    parser.add_argument('command', choices=['process', 'stream'])
    
    if command == 'process':
        parser.add_argument('--input-file', required=True)
        # Configure for file processing
    elif command == 'stream':
        parser.add_argument('--source', required=True)
        # Configure for network streaming
\end{verbatim}

While the current implementation treats production mode similarly to development mode (primarily differing in configuration defaults), the architectural separation facilitates future enhancements: production mode could integrate with telescope scheduling systems, implement real-time monitoring, or enforce access controls unavailable in development mode.

\subsection{Interactive Shell Mode}

An interactive mode complements command-line execution for exploratory workflows:

\begin{verbatim}
def interactive_shell():
    config = CorrelatorConfig()  # Start with defaults
    
    while True:
        print("\nCorrelator Interactive Shell")
        print("1. Configure parameters")
        print("2. Run correlation")
        print("3. Load configuration file")
        print("4. Show current configuration")
        print("5. Exit")
        
        choice = input("Select option: ")
        
        if choice == '1':
            # Prompt for parameter changes
            n_ants = input(f"Antennas [{config.n_ants}]: ")
            if n_ants:
                config.n_ants = int(n_ants)
            # ... repeat for key parameters
        
        elif choice == '2':
            run_correlator(config)
        
        elif choice == '3':
            path = input("Configuration file path: ")
            config = CorrelatorConfig.from_yaml(path)
        
        elif choice == '4':
            print(asdict(config))
        
        elif choice == '5':
            break
\end{verbatim}

This simple menu-driven interface suits initial system exploration and teaching contexts. Users adjust parameters through prompts, run correlations, and iterate without memorising command-line syntax. The interactive mode demonstrates configuration flexibility: the same \texttt{CorrelatorConfig} object used by command-line and programmatic interfaces drives the interactive shell.

\subsection{Batch Processing Script Interface}

For survey observations requiring repeated correlation with varying parameters, a batch script mode processes multiple configurations sequentially:

\begin{verbatim}
# batch_process.py
configs = [
    'config_observation1.yaml',
    'config_observation2.yaml',
    'config_observation3.yaml',
]

for config_file in configs:
    print(f"Processing {config_file}...")
    config = CorrelatorConfig.from_yaml(config_file)
    run_correlator(config)
    print(f"Completed {config_file}\n")
\end{verbatim}

This Python script imports the correlator as a library and invokes \texttt{run\_correlator} programmatically. The pattern enables sophisticated automation: scripts could generate configurations algorithmically, process results between correlations, or implement parallelism by distributing configurations across multiple processes or machines.

\subsection{Error Reporting and User Feedback}

The interface includes error handling that produces actionable messages. Configuration validation errors surface immediately:

\begin{verbatim}
$ correlator dev run --n-channels 300
Error: n_channels must be a power of 2, got 300
Valid values: 64, 128, 256, 512, 1024, 2048, 4096
\end{verbatim}

Runtime errors indicate the failure point and suggest remediation:

\begin{verbatim}
$ correlator dev run --input-file missing.npy
Error: Input file not found: missing.npy
Check file path and ensure file exists
\end{verbatim}

Successful completion reports key metrics:

\begin{verbatim}
Initializing FX correlator with 4 antennas, 256 channels
Creating simulated stream: 10.0s @ 1024.0 Hz
Processing...
  Saved integration 1 -> visibility_0001.npy
  Saved integration 2 -> visibility_0002.npy
  ...
  Saved integration 10 -> visibility_0010.npy
Complete: 10 integrations written to /workspace/outputs/
Configuration saved: /workspace/outputs/config.yaml
\end{verbatim}

These messages communicate progress without overwhelming users with implementation details. The balance between informativeness and verbosity supports both experienced users who need confirmation and novices who benefit from explicit feedback about system state.

In summary, the integration layer transforms the modular components into accessible, usable software through careful interface design, flexible configuration management, and clear user communication. The resulting system accommodates diverse usage patterns—command-line automation, interactive exploration, programmatic scripting—while maintaining a consistent underlying implementation.
